%% Appendix of the ARR submission. Built from the former supplementary material.
%% Included from main.tex after \appendix; ARR requires appendices in the same PDF,
%% after the references, in double-column format.

%% ============================================================
\section{Implementation of the Harness}

\subsection{Isolated scoring}
We evaluate every kernel in its own subprocess. The subprocess loads the generated module, runs the
correctness check and the timing loop, and writes the result as one JSON line on standard output.
The parent process reads that line and collects the results. This design is required, not optional.
The CUDA runtime keeps state per process, so if one bad kernel corrupts the context, then every
later evaluation in the same process fails as well. Without isolation, a sequential run over 250
tasks is not meaningful.

\subsection{Correctness check and chunked comparison}
We draw two random inputs and accept the kernel when the absolute error against the reference stays
below $10^{-2}$ on both (the framework's default path draws five, see Appendix~\ref{sec:app-variance}). A simple comparison allocates temporary tensors for the absolute difference, for
the tolerance, and for the mask. Each of these has the size of the output, so the peak memory grows
to three or four times the size of the output. On a 48GB GPU these temporary tensors alone made
\emph{correct} kernels fail with out of memory. To avoid this we compare in chunks of 64 million
elements and we also collect diagnostic statistics during the comparison. We free the input tensors
right after the forward pass, and we run the workers with an allocator option that reduces memory
fragmentation. For speed we take the median of ten runs, and the speedup is the reference time
divided by the generated time.

\subsection{Prompt structure}
The system prompt states the hardware facts of the target GPU. It gives the shared memory size,
which is about 100KB or 101{,}376 bytes, the amount of device memory, the memory bandwidth, and the
peak throughput. It also gives simple tile size rules, for example a block size of 64 or less for
matrix multiplication and a block size of 1024 for element-wise work. We also inject a small
database of Triton patterns and a set of architecture hints. The user prompt contains the full
PyTorch reference and the \textbf{input dtype contract}. The contract tells the model that it must
not guess the dtype of the reference inputs and must not cast them. Without this contract the model
writes a kernel that assumes fp16, and every such kernel fails on an fp32 reference. The repair
prompt adds the previous code, the error message, and the index of the current round. We require the
answer to be a single Python code block that contains a \texttt{Model} class derived from
\texttt{torch.nn.Module} and a kernel function with the \texttt{@triton.jit} decorator.

\subsection{Safety rules in the repair loop}
The loop enters the optimization stage only after the kernel is correct. If the values do not match,
we do not try to make the kernel faster. If the loop makes no progress for a fixed number of rounds,
we throw away the accumulated context and the previous code, and we generate again from scratch.
This second rule helps when the context is polluted, for example after an illegal memory access, but
it does not help when the model simply cannot solve the task. In our runs the two cases were easy to
tell apart. Tasks that a context reset rescued were different from tasks that failed even after two
resets, and the second group was concentrated on three-dimensional and four-dimensional tensor
products and on transposed matrix multiplication.

\subsection{Response parsing and multiple APIs}
The parser extracts the Markdown code block, renames the class to \texttt{Model}, and removes any
prose that is not Triton code. This step absorbs the different output styles of different models.
Commercial models use the chat completion API. Reasoning models do not accept a temperature
parameter, so they use the response API instead. Open models use the completion API of a local vLLM
server, and the quantized model uses llama.cpp. The client detects the right API from the model
identifier, falls back when the first call fails, and caches the result so that it does not repeat
the same detection on every call.

%% ============================================================
\section{Budget Settings and Wall-Clock Cost}\label{sec:app-budget}

The small budget is one sample and up to three repair rounds. We use it for full sweeps over 250
tasks. The large budget is three samples and up to nine repair rounds, which is at most 27 attempts
per task. The multi-turn curve changes only the number of effective turns, from 1 up to 14.

A full sweep under the large budget is dominated by wall-clock cost. On a slow local server, 27
attempts per task did not fit into a two-hour limit. As a result 59 percent of the tasks did not
finish, and the whole sweep took 22.9 hours. If we break down the result, 29 tasks were correct and
9 tasks were wrong after a normal evaluation, so only 38 tasks finished. The rest were 59 worker
timeouts and 3 tasks with no result. One point is worth noting. Several tasks first looked like
tasks that failed all 27 attempts, but in fact they were all timeouts and not wrong answers. A
reader who only looks at the surface accuracy can easily read this as a limit of model ability.

We draw three design rules from this. First, set the worker timeout from the number of attempts
times the average time per attempt, and not from a fixed guess. Second, keep the number of workers
below the number of concurrent serving slots. Third, draw the samples of best-of-$N$ in parallel and
not one after another.

Parallelism at the task level was harmful. When we ran two tasks at the same time, the serving slots
saturated, the throughput per slot dropped by more than half, generation exceeded the client
timeout, and the retry loop repeated forever, so no task finished at all. Sequential execution was
the better choice in our setting.

The repeated runs of Appendix~\ref{sec:app-variance} took 4.3 to 4.9 hours each on one RTX PRO 6000
that served the model and scored the kernels at the same time, about 32 GPU-hours in total. Every
full Level~1 sweep under the small budget has this cost on this GPU, and the Level~2 and Level~3
sweeps take longer because the reference models are larger.

%% ============================================================
\section{Raw Data of the Repair-Round Ablation}

Table~\ref{tab:ablation} gives the full numbers behind Table~\ref{tab:rq4a} of the main text.

\begin{table}[h]
\centering
\small
\setlength{\tabcolsep}{5pt}
\caption{Effect of the number of repair rounds on accuracy and on the error profile.
Model is gpt-4o and the task set is KernelBench Level~1.}
\label{tab:ablation}
\begin{tabular}{@{}lrrr@{}}
\toprule
\headrow{Rounds} & \headrow{Solved (/100)} & \headrow{Compile err.} & \headrow{Value mism.} \\
\midrule
1  & 19 & 59.3\% & 24.7\% \\
3  & 21 & 51.9\% & 39.2\% \\
5  & 19 & 45.7\% & 37.0\% \\
10 & 22 & 32.1\% & 53.8\% \\
\midrule
\textbf{Union} & \textbf{28} & --- & --- \\
\bottomrule
\end{tabular}
\end{table}

As the number of rounds grows, the share of compilation errors falls from 59.3 percent to 32.1
percent, while the share of value mismatches rises from 24.7 percent to 53.8 percent. Iterative
repair fixes syntax errors well, but it does not fix semantic errors on its own. This matches the
error transition finding of the main text. The union over the runs grows as 19, 25, 26, and 28, so
it always stays above the best single run, which solved 22 tasks.

%% ============================================================
\section{Details of Model Cascading}

The main text only reports the conclusion of the cascade experiment. Table~\ref{tab:cascading}
gives the accumulated result of each stage.

\begin{table}[h]
\centering
\small
\setlength{\tabcolsep}{4pt}
\caption{Accumulated accuracy of the cascade. Each stage only runs on the tasks that the previous
stage failed.}
\label{tab:cascading}
\begin{tabular}{@{}lrrrr@{}}
\toprule
\headrow{Configuration} & \headrow{L1} & \headrow{L2} & \headrow{L3} & \headrow{Total (/250)} \\
\midrule
gpt-4o-mini            & 15 & 1  & 0 & 16 \ (6.4\%) \\
gpt-4o                 & 19 & 6  & 0 & 25 \ (10.0\%) \\
gpt-4o $+$ gpt-5.2     & 42 & 15 & 4 & 61 \ (24.4\%) \\
\quad $+$ gpt-5.2-pro  & 43 & 15 & 4 & 62 \ (24.8\%) \\
\midrule
claude-opus-4-7 alone  & 46 & -- & -- & -- \\
opus $+$ gpt-5.4-nano  & 48 & -- & -- & -- \\
union of five models   & \textbf{51} & -- & -- & -- \\
\bottomrule
\end{tabular}
\end{table}

We observe three things. First, adding the next generation model to gpt-4o raises the total from 25
tasks to 61 tasks, which is a factor of 2.44. The 36 new tasks are 23 tasks on Level~1, 9 tasks on
Level~2, and 4 tasks on Level~3. Second, the expensive reasoning model adds almost nothing. It ran
on 225 tasks and added only one task, a Level~1 task with a speedup of 7.60 times. If we also
consider its cost and its response time, its value inside a cascade is close to zero. Third, one
recent flagship model alone, with 46 tasks, already beats the earlier three-model cascade with 43
tasks. The union of five models, with 51 tasks, still beats the best single model by 10.9 percent,
so the value of independent samples does not disappear in the newest generation.

%% ============================================================
\section{Error Distribution by Level}

Table~\ref{tab:rq3} of the main text only covers Level~1. Table~\ref{tab:error_detail} covers all three levels.
The timeout column counts our own response and evaluation time limits together with unclassified
runtime errors. It is not a constraint of the benchmark.

\begin{table}[h]
\centering
\small
\setlength{\tabcolsep}{4pt}
\caption{Share of each error type among failed tasks, by model and by level.}
\label{tab:error_detail}
\resizebox{\columnwidth}{!}{%
\begin{tabular}{@{}llrrrrr@{}}
\toprule
\headrow{Model} & \headrow{Lv.} & \headrow{Failed} & \headrow{Compile} & \headrow{Value} & \headrow{CUDA} & \headrow{Timeout} \\
\midrule
\multirow{3}{*}{gpt-4o-mini}
  & L1 & 85 & 71.8\% & 20.0\% & 2.4\% & 2.4\% \\
  & L2 & 99 & 70.7\% & 22.2\% & 4.0\% & 1.0\% \\
  & L3 & 50 & 82.0\% & 6.0\%  & 6.0\% & 0\% \\
\midrule
\multirow{3}{*}{gpt-4o}
  & L1 & 81 & 45.7\% & 37.0\% & 11.1\% & 0\% \\
  & L2 & 94 & 51.1\% & 41.5\% & 4.3\%  & 0\% \\
  & L3 & 50 & 38.0\% & 56.0\% & 4.0\%  & 0\% \\
\midrule
\multirow{3}{*}{gpt-5.2}
  & L1 & 58 & 6.9\%  & 58.6\% & 1.7\% & 20.7\% \\
  & L2 & 85 & 25.9\% & 62.4\% & 0\%   & 8.2\%  \\
  & L3 & 46 & 17.4\% & 71.7\% & 0\%   & 4.3\%  \\
\midrule
gpt-5.5         & L1 & 63 & 71\% & 3\%  & 0\% & 25\% \\
claude-opus-4-7 & L1 & 54 & 0\%  & 65\% & 0\% & 28\% \\
\bottomrule
\end{tabular}}
\end{table}

The error transition appears at every level. For gpt-4o-mini the compilation error is dominant at
all three levels, and for gpt-5.2 the value mismatch is dominant at all three levels. Inside one
model the same shift also appears with difficulty. For gpt-4o the share of value mismatches grows
from 37.0 percent on Level~1 to 56.0 percent on Level~3, so failures become more semantic as the
task becomes harder.

For the reasoning model gpt-5.2-pro, 100 percent of the Level~1 failures and 81.9 percent of the
Level~2 failures were timeouts. On Level~3 all failures came from API rate limits. That is a
different cause, so we do not merge it into the timeout column.

%% ============================================================
\section{Details of the Speedup Results}

We only measure speedup on kernels that pass the correctness check. Table~\ref{tab:fast} reports the
result.

\begin{table}[h]
\centering
\small
\setlength{\tabcolsep}{6pt}
\caption{Share of correct kernels that reach the Fast-1 and Fast-2 thresholds. The denominator is
the number of correct kernels for that model.}
\label{tab:fast}
\begin{tabular}{@{}lrr@{}}
\toprule
\headrow{Model} & \headrow{Fast-1} & \headrow{Fast-2} \\
\midrule
gpt-4o-mini     & 7/16 \ (43.8\%)  & 1/16 \ (6.3\%)  \\
gpt-4o          & 12/25 (48.0\%)   & 4/25 \ (16.0\%) \\
gpt-5.2         & 25/36 (69.4\%)   & 5/36 \ (13.9\%) \\
gpt-5.2-pro     & 10/13 (76.9\%)   & 3/13 \ (23.1\%) \\
gpt-5.4-nano    & 21/35 (60.0\%)   & 10/35 (28.6\%)  \\
gpt-5.5         & 22/37 (59.5\%)   & 8/37 \ (21.6\%) \\
claude-opus-4-7 & 29/46 (63.0\%)   & 14/46 (30.4\%)  \\
\bottomrule
\end{tabular}
\end{table}

The Fast-2 share grows from 6.3 percent to 30.4 percent across model generations, so newer models
write not only more correct kernels but also better optimized ones. The reasoning model is an
interesting exception. It solved only 13 tasks, which is the smallest number in the table, yet its
Fast-2 share is 23.1 percent, which is high. When a reasoning model does solve a task, it tends to
write a well optimized kernel. Its overall success rate is low, however, so its practical value
stays limited.

The largest single speedup that we observed was 84.03 times on one Level~3 task. Other high cases
were 25.45 times, 19.10 times, and 11.86 times. These cases show that a generated Triton kernel can
beat the default PyTorch implementation by a wide margin on some operation patterns. Speedup still
depends on the hardware baseline, so these numbers are only comparable inside one hardware setting.

%% ============================================================
\section{Details of the Reused Open Artifacts}

The main text only reports the result of each port. Here we describe how each artifact is built and
what happened when we ran it.

\subsection{Failure-rule memory}
The system turns each failure into one short rule of the form ``you must not do this'', stores the
rule together with a score, removes duplicates after normalization, and injects the top rules into
the generation prompt. In one full evaluation the store collected between 73 and 80 distinct rules.
Examples are a rule that says the kernel body must not call PyTorch operations, rules about the
dimensions of a matrix multiplication and about the block size along the reduction axis, and rules
about contiguity assumptions. The cost is one extra model call per failure, which is cheap, and this
was the most portable artifact that we reused.

The risk, however, is not symmetric. A good rule saves a few attempts, while a wrong rule with a
high score can block a whole class of tasks. In our runs the rules that the system collected while
the evaluation harness was still broken polluted the store. The polluted store held 320 rules. Some
of them treated an infrastructure defect as a rule for the model, for example a rule that said the
GPU does not have enough memory to run the kernel. These rules had a high score and filled the top
of the prompt. As a result the model failed on tasks as simple as a square matrix multiplication,
even with nine repair rounds and three samples. Accuracy recovered after we cleared the store. There
are two lessons. The store must be cleared after any infrastructure fix, and in a multi-server setup
with a shared store the pollution spreads from one server to the other.

\subsection{The Triton-specific reinforcement learning model}
We served the released weights on a local vLLM server. This model does not support the system role
of the chat format, so our first calls failed. We changed the client to fold the system message into
the user prompt, and the calls then worked. On the first 20 tasks of Level~1 the model reached 0.85
on the 96GB GPU and 0.95 on the H100. The H100 result is higher because the model needs only about
16GB, so the scorer still had more than 55GB and no task failed with out of memory. Under the same
conditions the 30B general open model reached only 0.60 on the H100. A small specialized model
therefore beats a large general model and at the same time leaves more memory for the scorer, which
matters in practice.

\subsection{The open deep-agent framework}
We cloned the released framework and ran it with our own open models. Three problems appeared during
the port. First, the serving context length decides whether repair works at all. With a 16K context
the repair prompt, which grows with the code and the error text, exceeded the input and output
limit, and every round was lost. With a 64K context the same tasks were solved in a few tens of
seconds. Second, installing an external project into a shared environment breaks the dependencies.
The editable install of the framework raised the version of shared packages, and the serving engine
then failed to start. An external project must be installed into its own virtual environment. Third,
the shell command returned an exit code of zero even when the HTTP response was an error, and we
twice concluded that the server was healthy when it was not. A health check must look at the status
code of a real completion call.

We also note that this framework verifies its answers with tests that the model itself writes, and
its tolerance is looser than ours. Its accuracy numbers and our accuracy numbers therefore use
different verification strengths, which is one more reason why numbers from different systems should
not be compared directly.

\subsection{The multi-turn reinforcement learning environment}
We reproduced the released environment on the 100 tasks of Level~2. Table~\ref{tab:repro} gives the
metrics that the main text does not list.

\begin{table}[h]
\centering
\small
\setlength{\tabcolsep}{7pt}
\caption{Reproduction of the released multi-turn reinforcement learning environment on KernelBench
Level~2 with 100 tasks. All values are computed over all tasks.}
\label{tab:repro}
\begin{tabular}{@{}lrr@{}}
\toprule
\headrow{Metric} & \headrow{Original paper} & \headrow{Our reproduction} \\
\midrule
Fast@1   & 49.2\% & 42.9\% \\
Fast@1.2 & 25.6\% & 24.3\% \\
Fast@1.5 & 7.4\%  & 12.9\% \\
Fast@2   & 2.1\%  & 7.5\%  \\
\bottomrule
\end{tabular}
\end{table}

Fast@1.2 is almost the same and Fast@1.5 and Fast@2 are higher, so we judge the reproduction to be
successful. Fast@1 is 6 points lower, which we attribute to a different GPU and to sampling noise.
Correctness alone was 76.3 percent, and pass@1 under a loose threshold was 85.7 percent.

The infrastructure cost was high. The grading path of the trainer needs the FSDP actor and the vLLM
rollout engine in the same memory, and five attempts on a single 96GB GPU all failed. An eight-GPU
node was in practice the minimum requirement. We also had to clear the existing cluster completely
before each run, because leftover processes of a different version attached themselves to our job.
The released validation data covers Level~2 only, so the same reproduction on Level~1 and Level~3
would need a new data set in the same format. The point of this subsection is that reproducibility
does not imply accessibility.

%% ============================================================
\section{Details of the Quantization Study}

We loaded a very large open reasoning model with aggressive quantization and evaluated it under the
same protocol. We served it with layer-split pipeline parallelism and we kept one separate GPU for
scoring.

The 2-bit version was about 254GB. It reached 0.10 on Level~1 and solved nothing at all on Level~2,
on Level~3, and on the real serving kernels. It needed about 8.5 hours per level and 21.5 hours in
total. It spent about 2{,}600 tokens of thinking per attempt with no gain in accuracy. A general
reasoning model with two orders of magnitude more parameters therefore stayed below half of the 8B
specialized model.

The 4-bit version was about 467GB, which was the highest precision that fit into our serving budget.
Accuracy recovered to 0.96 over the tasks that finished inside the time limit, which confirms that
quantization loss was the cause of the collapse. This version, however, generated only about 8
tokens per second. As a result 76 percent of the tasks exceeded the two-hour limit, the surface
accuracy was only 0.23, and the run took 30 hours.

Taken together, under a fixed GPU budget a very large model at low precision loses quality, and the
same model at full precision loses speed. A mid-size model at full precision is the better choice
for kernel generation.

%% ============================================================
\section{Details of the Cross-Benchmark Study}

\subsection{The Triton-specific benchmark}
This benchmark has 166 tasks and uses a different evaluation style from KernelBench. Each task is a
PyTorch function together with a test harness, so a scorer that expects a module and an input
generator cannot run it. We therefore wrote a separate runner. The runner executes the reference and
the generated code with the same test harness, compares the two result dictionaries, and reports
both whether the code runs and whether the results match. Over the four combinations of two models
and two GPUs, execution accuracy stayed inside a narrow band between 0.151 and 0.187, and the
advantage of the specialized model that we saw on KernelBench disappeared. The suite also contains
tasks such as optimizer updates, which do not map naturally onto a Triton kernel, so a part of the
suite is structurally out of reach.

To find the cause of the gap against the original paper, we ran a two-by-two ablation over the prompt
format and the budget. Aligning the prompt with the training format of the specialized model alone
did not help, and execution accuracy even fell from 0.151 to 0.127. Turning on the budget raised
execution accuracy by about 2.3 times and call accuracy by about 2.5 times. Under the larger budget
the conversion-style prompt, which provides the PyTorch reference, was in fact better than the
aligned prompt. The gap is therefore explained by single-attempt evaluation and not by a mismatch
with the training prompt.

\subsection{The standardized agent environment}
This environment has 139 standard tasks and 13 out-of-distribution tasks. Under a single attempt,
106 of the 139 standard tasks compiled, and 17 tasks passed both compilation and the correctness
check, which is 12.2 percent. The geometric mean of the performance of the passing kernels relative
to the oracle was 0.143, so they were about seven times slower than the oracle, and only three
kernels were faster than the oracle. The main cause of compilation failure was an import error, with
29 cases in which the model used a module or a symbol that does not exist. On the 13
out-of-distribution tasks, 12 tasks compiled but only one task passed, which is 7.7 percent.

With a test-time scaling agent, Pass@1 on the standard tasks rose to 18.0 percent and the geometric
mean of performance rose to 0.285, which is about twice as high. The compilation rate, however, fell
to 45.3 percent, because the last proposal of a repair loop sometimes does not compile. On the
out-of-distribution tasks Pass@1 stayed at 7.7 percent and did not move at all. This is the
measurement behind the boundary result of the main text.

\subsection{The track for a new domain-specific language}
The same environment has a track with 12 tasks for a new domain-specific language. Nine of these
tasks could not run in our environment at all. Some of them need a separate compiler fork, and
others use an API that only exists in a newer version. On the three tasks for which the oracle could
run, the specialized reinforcement learning model passed none under a single attempt. The sample is
small, so we do not draw a quantitative conclusion. We report it as evidence that the collapse
outside the training distribution is not limited to one axis of one benchmark.

\subsection{The real serving kernel benchmark}
We took 8 dense tasks from a benchmark of real serving kernels for a recent NVIDIA architecture and
converted them into the KernelBench format. The conversion turns the reference definition and the
tensor specification into a module and an input generator. Among the four combinations, only the
specialized model on the H100 reached 0.25, and the other three stayed at zero. This points to
differences in the Triton version and in the compute architecture. The attention and mixture-of-expert
tasks need structured input synthesis, so we leave them for future work.

%% ============================================================
\section{Beyond the GPU: Two NPU Targets}\label{sec:app-npu}

To test whether the findings are properties of GPUs and Triton or of the generate-and-verify loop
itself, we ported the harness, unchanged in its structure, to two neural processing units with
different programming models. The first exposes a Triton dialect, so the language is familiar to the
models. The second exposes a Rust-embedded domain-specific language that no model has seen at scale.
Both experiments use the same A/B design as the rule-memory study: condition A gives the model the
vendor documentation alone, and condition B adds one page of constraints that we measured on the
device.

\subsection{A Triton dialect on a commercial NPU}\label{sec:app-npu-triton}

The first target is a data-center NPU whose SDK accepts Triton kernels through a vendor compiler.
We took five operators that a transformer block needs, namely GELU, layer normalization, matrix
multiplication, attention, and a block-scaled product, and asked six current gateway models for
three samples each, which gives 90 candidates per condition. Every candidate passes the compile gate
on the device host and is then executed on the NPU against the PyTorch reference, with the
signal-to-quantization-noise ratio as the correctness measure and the median of three latency runs.

\begin{table}[t]
\centering
\small
\setlength{\tabcolsep}{6pt}
\caption{Correct kernels out of 15 (five operators, three samples) on the Triton-dialect NPU.
Condition A gives the vendor documentation; condition B adds one page of measured constraints.
No repair round was ever triggered, because every failing kernel compiled.}
\label{tab:app-npu-triton}
\begin{tabular}{@{}lrr@{}}
\toprule
\headrow{Model} & \headrow{A} & \headrow{B} \\
\midrule
claude-opus-4-8        & 12/15 & 15/15 \\
claude-sonnet-5        & 12/15 & 15/15 \\
gpt-5.5                & 12/15 & 15/15 \\
gemini-3.1-pro-preview & 12/15 & 15/15 \\
deepseek-v4-pro        & 10/15 & 15/15 \\
qwen3.8-max            &  8/15 & 13/15 \\
\midrule
all models             & 66/90 & 88/90 \\
\quad of which GELU    &  0/18 & 18/18 \\
\bottomrule
\end{tabular}
\end{table}

Table~\ref{tab:app-npu-triton} shows the result. Under condition A, 85 of the 90 candidates
compile but only 66 are correct, and the gap has one dominant cause. Sixteen of the 24 failures are
kernels that compile, run, and return uninitialized memory. The vendor's supported-operation list
includes the sigmoid, the error function, the minimum, and the conditional select, but none of them
lowers on the device, and a Python scalar literal inside a tensor expression does not lower either.
When lowering fails, the runtime silently falls back and the output buffer is never written. GELU
needs exactly these operations, so all 18 GELU candidates compile and all 18 are wrong, for every
model. The compile gate is blind to this failure, and so is any report that counts compiled kernels.
Only execution with a numerical comparison catches it, which is the verification-strength axis of
Finding~2 on a second platform.

Condition B adds one page that lists which operations really lower, how to build a constant from a
tensor, and how to compose a sigmoid from the exponential. With this page the five flagship models
reach 15 of 15 and the total rises from 66 to 88 of 90; the two remaining failures are gateway
timeouts. The generated kernels run at the latency of the hand-written reference, with a median ratio
of 0.99 in both conditions. A local 8B open model under the same setup solved only the matrix
multiplication, one of three operators we tried, so the result above is a property of the flagship
generation. Measured platform knowledge, then, transfers completely when the model already writes
the language fluently, and the number that moved was not the compile rate but the share of compiled
kernels that are correct.

\subsection{A Rust-embedded DSL on a second NPU}\label{sec:app-npu-dsl}

The second target shares nothing with Triton: three attention and feed-forward kernels of a 12B
decoder, written in a Rust-embedded domain-specific language for a different NPU, taken from a public
kernel-optimization challenge. The harness adds a contract gate that rejects edits outside the
permitted files or to any device-function signature, and the compile gate classifies the result into
a Rust front-end error, a backend error, or a compiling kernel with a static schedule length.
Because a compiling kernel is not an improved kernel, two no-op detectors follow the compile gate:
identity of the source after comment removal, and identity of the static schedule after source
positions are stripped. We ran the 30B open coder with three samples and three repair rounds per
kernel under both conditions.

\begin{table}[t]
\centering
\small
\setlength{\tabcolsep}{5pt}
\caption{The 30B open coder on three NPU kernels in a Rust-embedded DSL, small budget. A is the
vendor documentation and B adds the measured constraints. A candidate counts as an improvement only
if it compiles, changes the static schedule, and shortens it. The last column counts contract
violations.}
\label{tab:app-npu-dsl}
\begin{tabular}{@{}lrrrr@{}}
\toprule
\headrow{Cond.} & \headrow{Compiled} & \headrow{No-op} & \headrow{Improved} & \headrow{Violations} \\
\midrule
A & 7/9 & 7 & 0 & 1 \\
B & 7/9 & 6 & 0 & 3 \\
\bottomrule
\end{tabular}
\end{table}

Table~\ref{tab:app-npu-dsl} shows the result over 42 attempts and about 400K output tokens.
Fourteen of eighteen samples compiled, and thirteen of the fourteen were no-ops: the model added
comments, renamed a helper, or undid its own edit in the repair round to satisfy the compiler. The
one candidate that changed the schedule left its length unchanged. Over all attempts the Rust front
end rejected 57 percent, and about four fifths of those rejections were one grammar rule, a mapping
macro used in a type position. No candidate reached the backend compiler. Here the measured
constraints did not help. The page describes what to do, the model could not express it in the
language, and the contract violations rose from one to three. The hand-optimized kernels on the same
tasks run about seven times faster than the baseline, and the model recovered none of that gap.

The two targets together sharpen Finding~8. The same page of measured knowledge moved accuracy from
66 to 88 of 90 where the model already spoke the language and moved nothing where it did not. Budget
and knowledge act after the syntax wall, not before it, and the stage at which a target fails decides
what any further investment can buy.

%% ============================================================
\section{Details of the Cross-Hardware Study}

Under the same conditions the 30B open model reached 0.75 on the 96GB GPU and 0.60 on the H100 over
the first 20 tasks of Level~1. The main cause is GPU memory. The H100 has 80GB, so after serving the
30B model the scorer had only about 43GB, while the 96GB GPU left about 52GB. For tasks with large
inputs this difference alone decides whether the task fails with out of memory. If we correct for
the out-of-memory cases, the H100 number rises to about 0.70, which is close to the 96GB GPU. We
also tried to lower the serving share further, but then the server did not start at all because the
key-value cache became too small, so this environment has a lower bound.

Speedup behaves in the opposite way and depends strongly on the hardware. On a square matrix
multiplication we measured 1.61 times on the 96GB GPU and 0.83 times on the H100. The cuBLAS
baseline of the H100 is very strong, so a generated Triton kernel is at a disadvantage there. The
median speedup was 1.51 against 1.44, and the Fast-2 share was 0.10 against 0.05.

We repeated the run on a third GPU with 48GB and obtained a similar profile across levels, with 0.37
on Level~1, 0.12 on Level~2, and 0.04 on Level~3. This supports the conclusion that accuracy is
almost independent of the GPU when memory is sufficient. This GPU, however, produced only about one
fourth to one fifth of the throughput of the H100, so the same sweep took about three times longer
in wall-clock time. Accuracy is independent of the hardware, but whether the experiment can be run
at all is not.

%% ============================================================
\section{Practical Notes for Reproduction}\label{sec:notes}

Several operational factors distorted our measurements more than once, and we list them here
because they are easy to miss and hard to detect from the accuracy number alone.

First, serving and scoring must not compete for the same memory. If the serving process takes too
large a share of the GPU, even a trivial kernel fails, as Finding~2 shows. Second, a released model
artifact can carry its own generation limit, and the server clamps the client limit to that value,
so raising the limit on the client side alone does not remove truncation. Third, a reasoning model
spends thousands of tokens before it writes code, so the common response limit of the agent
framework is too small and cuts the answer before the code block. Fourth, the serving context length
must hold the repair prompt, which grows with the code and the error text, otherwise every repair
round is lost. Fifth, an experience memory that was filled while the harness was still broken keeps
poisoning later runs, so the store must be cleared after an infrastructure fix. Sixth, a shared
store spreads that damage across servers, so each experiment should use its own path.

We also found that parallelism at the task level can be harmful. When we ran two tasks at the same
time, the serving slots saturated, the throughput per slot dropped, generation exceeded the client
timeout, and the retry loop never terminated, so no task finished at all. Sequential execution was
the better choice on our hardware.

%% ============================================================
\section{Run-to-Run Variance and Harness Sensitivity}\label{sec:app-variance}

All numbers in the main text are single runs except the Level~1 accuracy of the 8B open model
under the small budget, which we repeated to measure how much one run of this benchmark moves.
Table~\ref{tab:app-variance} lists the runs. The sampling temperature is 0.2, so the variance comes
from sampling in generation and repair, not from the scorer, which is deterministic for a given
kernel.

\begin{table}[t]
\centering
\small
\setlength{\tabcolsep}{5pt}
\caption{Repeated runs of the 8B open model on all 100 tasks of Level~1 under the small budget.
The first row uses the harness of the main text (runs: 0.47, 0.44, 0.39, 0.47) and the second row
uses the framework's own command-line path (runs: 0.31, 0.41, 0.33). Union and intersection count
the tasks solved by any run and by every run of a row.}
\label{tab:app-variance}
\begin{tabular}{@{}lrrrr@{}}
\toprule
\headrow{Path} & \headrow{Runs} & \headrow{Mean $\pm$ sd} & \headrow{Union} & \headrow{Inter.} \\
\midrule
main-text harness & 4 & 0.44 $\pm$ 0.04 & 59 & 25 \\
command-line path & 3 & 0.35 $\pm$ 0.05 & 50 & 23 \\
\bottomrule
\end{tabular}
\end{table}

Two observations follow. First, the standard deviation of a single full sweep is four to five
points. Any difference of that size between two rows of a table in this paper is within noise, and
only the gaps between model generations, which are ten points and more, are safe to interpret.
Second, the set of solved tasks moves more than the count. The four runs of the main-text harness
agree on only 25 tasks, yet together they cover 59, which is twelve more than the best single run.
This is the same effect that Finding~4 reports for repair depth, and it appears here without any
change of budget. Breadth, in the form of independent runs, buys coverage that depth does not.

The second block is a harness effect rather than a model effect. Both paths load the same model,
the same prompt templates, and the same task set with one sample and three repair rounds. They
differ in five settings. The main-text harness runs one operating-system process per task in
sequence, waits up to 2400 seconds per task, samples at temperature 0.2, verifies on two random
inputs, and keeps profiling off. The command-line path runs two tasks in parallel threads, stops
each evaluation after 120 seconds, samples at temperature 0.3, verifies on five random inputs, and
runs the profiler. None of these settings is part of the model, and together they move the mean by
nine points, which is more than twice the run-to-run standard deviation. This is why we report the
settings of Section~\ref{sec:design} with every number and why Finding~2 lists them among the axes
that an accuracy report must state.

%% ============================================================
\section{Hidden Variables and a Reproduction Checklist}

We list the factors that distorted our measurements more than once. They have one thing in common.
None of them appears in an accuracy report.

\begin{itemize}\itemsep3pt
\item \textbf{Keep serving memory and scoring memory apart.}
If the serving process takes too large a share of the GPU, even a trivial kernel fails with out of
memory. When serving and scoring share one GPU, the serving share must match what the scorer needs.
\item \textbf{Check the generation limit inside the model artifact.}
The maximum generation length in the released generation config is clamped by the server. Raising
the limit on the client side alone does not remove truncation, so the artifact itself must be
checked.
\item \textbf{Raise the response token limit of the framework.}
A reasoning model spends thousands of tokens on its thinking trace before it writes any code. A
common limit of 8K cuts the answer before the code block and turns a solvable task into a task with
no code. In our runs, raising the limit to 13{,}000 moved accuracy on the affected range from 0.37
to 0.78.
\item \textbf{Make the serving context long enough.}
The repair prompt grows with the code and with the error text. If the context is short, every repair
round is lost.
\item \textbf{Clear the experience memory after an infrastructure fix.}
Rules that were collected while the harness was broken keep lowering accuracy in a healthy
environment.
\item \textbf{Give each experiment its own store path.}
When several servers share one experience store, the pollution spreads from one experiment to the
next.
\item \textbf{Prefer sequential execution on a small serving cluster.}
If the serving slots saturate, the throughput drops and the retry loop does not terminate.
\end{itemize}

Finally, we repeat the axes that should be reported next to any accuracy number. These are the
attempt budget, the verification strength, the memory headroom of the scorer, the serving context
length, the response token caps of both the framework and the model artifact, and the evaluation
distribution. A comparison of numbers without these axes is a comparison of measurements from
different regimes.

